\chapter{Algebraic Curves}
\section{Main Tools}

The main tool we will use in this chapter is the following, and will be
stated and proved for polynomials, even though it would be more general.

\begin{theorem}{Implicit function theorem}\label{thm:ift}
    Let \(f \in  \mathbb{C}[z, w]\). Consider \(\mathcal{C} = \{{(x, y)} \in
    \mathbb{C}^2 : f{(x, y)} = 0\} \). Assume then \({(z_{0}, w_{0})} \in
    \mathcal{C}\) and \(\frac{\partial f}{\partial w} {(z_{0}, w_{0})} \neq 0\).
    Then there exists \(U \ni z_{0}\) and \(V \ni w_{0}\) open sets and an
    holomorphic map \(g : U \to \mathbb{C}\) holomorphic such that
    \[
      \mathcal{C} \cap U \times  V = \{{(z, g{(z)})} : z \in U\} 
    \]
\end{theorem}

\begin{proof}{}
    Use first the \(C^{\infty}\)-implicit function theorem. Write then \(f {(z,
    w)} = \sum \alpha_{l,j} z^{l} w^{j} \) and we have 
    \[
      \partial_w f = \sum \alpha_{l, j} z^{l} j w^{j-1} dw = \frac{\partial
      f}{\partial w} dw \implies \partial_w f{(z_{0}, w_{0})} = \underbrace{\frac{\partial
      f}{\partial w}{(z_{0}, w_{0})}}_{\neq 0} dw
    \]
    hence \(\partial_w f{(z_{0}, w_{0})}\) is invertible and thus there exists
    \(U, V\) open sets and \(g : U \to \mathbb{C}\), \(g \in C^{\infty}\) such
    that \(\mathcal{C} \cap {(U \times  V)} = \{{(z, g{(z)})} : z \in U\} \).

    We just now need to show that \(g\) is holomorphic. Set \(F : U \to
    \mathbb{C}\) given by \(F{(z)} := f{(z, g{(z)})} = 0\) on \(U\). Then
    \[
      d_z F = 0 \iff \frac{\partial F}{\partial z} dz + \frac{\partial
      F}{\partial \overline{z}} \overline{z} = 0
    \]
    Hence
    \[
      \frac{\partial F}{\partial \overline{z}} = 0 = \frac{\partial f}{\partial
      w} {(z, g{(z)})} \cdot \frac{\partial g}{\partial \overline{z}}{(z)}
      \implies \frac{\partial g}{\partial \overline{z}} \equiv 0 \text{ near }
      z_{0}
    \]
    hence \(g\) is holomorphic.
\end{proof}

\begin{theorem}{}
    Let \(f \in \mathbb{C}[z,w]\) such that 
\begin{enumerate}[label = \arabic*.]
    \item \(\forall {(z, w)} \in \mathcal{C}\), \({\left( \frac{\partial
        f}{\partial z}, \frac{\partial f}{\partial w} \right)}\neq {(0, 0)}  \) 
    \item \(f\) is irreducible in \(\mathbb{C}[z,w]\) 
\end{enumerate}
Then there exists a complex structure on \(\mathcal{C}\) such that the two
projections \(p_{1}, p_{2} : \mathcal{C} \to \mathbb{C}\), given by \({(z, w)}
\mapsto z\) and \({(z, w)} \mapsto w\) respectively, are holomorphic.

\end{theorem}

\begin{remark}{}
    The two regularity requisistes ensure that (1.) we can use the implicit
    function theorem and (2.) that \(\mathcal{C}\) is connected. This second
    result would require a deeper explanation coming from algebraic geometry.
\end{remark}

\begin{proof}{}
    By \({(1.)}\) if \({(z_{0}, w_{0})} \in \mathcal{C}\) then one of the two
    derivatives is nonzero, let's assume \(\frac{\partial f}{\partial w} \neq
    0\). Using the IFT (theorem~\ref{thm:ift}) we know that 
    \[
        \underbrace{\mathcal{C} \cap {(U \times  V)}}_{\mathcal{O}_0}  = \{{(z, g{(z)})} : z \in U\} \text{ with
      \(g\) holomorphic}
    \]
    Then define \(\varphi_{0} : \mathcal{O}_0 \to \mathbb{C}\) by
    \(\varphi_{0}{(z, w)} = z = p_{1}{(z)}\). Then \({(\mathcal{O}_0,
    \varphi_{0})}\) is a chart. Similarly, if \(\frac{\partial f}{\partial z}
    \neq 0\) then the cart will be \({(\mathcal{O}_0, \varphi_{0} = p_{2})}\).

    We say that
    \[
      \bigcup_{{(z_{0}, w_{0})} \in  \mathcal{C}} {(\mathcal{O}_0, \varphi_{0})}
      \text{ is an atlas}
    \]
    indeed, if \(\mathcal{O} \cap \mathcal{O}' \neq \varnothing\), then
    \(\varphi' \circ \varphi ^{-1}{(z)} = g{(z)}\) holomorphic. %TODO exercise

    Finally, note that \(p_{1}\) and \(p_{2}\) are holomorphic because they are
    charts.
\end{proof}

\section{Projective Curves}

Let \(F \in  \mathbb{C}[x, y, z]\) such that \(F\) is homogenous, i.e. \(\forall
\lambda \in \mathbb{C}\), \(F{(\lambda x, \lambda y, \lambda z)} =
\lambda^{d}F{(x, y, z)}\) with \(d \in \mathbb{Z}_{\ge 0} \). Then we define
\[
    \mathcal{C} = \{[x, y, z] \in \mathbb{P}^2 \mathbb{C} : F{(x, y, z)} = 0\} 
\]
which makes sense because \(F\) is homogenous.

\begin{definition}{Smooth}
    A projective curve \(\mathcal{C}\) is said to be \textbf{smooth} if
    \(\forall [x, y, z] \in \mathcal{C}\), \(dF \neq 0\) 
\end{definition}

\begin{theorem}{}
    Assume \(\mathcal{C}\) is \emph{smooth}. Then there exists a complex
    structure on \(\mathcal{C}\) such that \(\mathcal{C}\) is a compact
    connected Riemann Surface. 
\end{theorem}

\begin{proof}{}
    v. \cite{miranda} p. 15
\end{proof}

\begin{eser}{}
    \(\mathcal{C} \subseteq \mathbb{P}^2 \mathbb{C} \) is smooth
    \(\iff\) \(\forall i=1,2,3\), \(U_{i} \cap \mathcal{C}\) is \emph{affine
    smooth}.
\end{eser}

\begin{example}{}
    Let \(F{(x, y, z)} = x^{d} + y^{d} + z^{d}\). Then \(F\) is homogenous of
    degree \(d\), hence it defines a projective curve \(\mathcal{C}\).
    Moreover, \(\mathcal{C}\) is smooth, since \(\frac{\partial F}{\partial
    \star} = d \star^{d-1}\), with \(\star = x, y, z\).
\end{example}

\begin{remark}{}
    If \(\mathcal{C} = \{f{(u, v)} = 0\} \subseteq \mathbb{C}^2 \) is affine,
    one can consider \(F{(x, y, z)} = z^{d} f{(\frac{x}{z}, \frac{y}{z})}\) and
    \(\widehat{\mathcal{C}} = \{[x, y, z] \in \mathbb{P}^2 \mathbb{C} : F{(x, y,
    z)} = 0\} \) and \(\widehat{\mathcal{C}} \cap U_{3} ``= \mathcal{C}''\). If
    \(\mathcal{C}\) is smooth, then \(\widehat{\mathcal{C}}\) is in general non
    smooth.
\end{remark}

\begin{eser}{}
    Consider \(\mathcal{C} = \{{(x,y)} \in \mathbb{C}^2 : y^2 = P{(x)}\} \)
    where \(P \in \mathbb{C}[X]\) non constant and \(P\) has simple roots.
\begin{enumerate}[label = \arabic*.]
    \item Check that \(P\) is smooth
    \item Check that \(\widehat{\mathcal{C}}\) is smooth \(\iff d \le 3\) 
\end{enumerate}

\end{eser}

